\section{Experiments}
\label{sec:experiments}

\subsection{Dataset}
We evaluate on the LineMOD dataset~\cite{hinterstoisser2012model}, a widely used benchmark for 6D pose estimation. It contains 13 objects with approximately 1,200 images per object captured under varying viewpoints. We use an 80/10/10 train/val/test split with interleaved sampling to ensure viewpoint diversity across splits.

\subsection{Implementation Details}
\textbf{Object Detection.} We train YOLOv8~\cite{yolov8} to detect objects and provide bounding boxes for the pose estimation network.

\textbf{Training.} We train for 150 epochs using AdamW optimizer~\cite{loshchilov2019adamw} with learning rate $10^{-4}$, weight decay $10^{-4}$, and batch size 32. We use cosine annealing learning rate schedule~\cite{loshchilov2017sgdr}. Loss weights are $\lambda_r = 2.0$ and $\lambda_t = 5.0$.

\textbf{Data Augmentation.} We experimented with color jitter and bounding box jitter augmentations but found they did not improve performance on LineMOD, likely due to the dataset's controlled indoor environment with consistent lighting.

\subsection{Main Results}

\cref{tab:main_results} shows our main results on LineMOD.

\begin{table}[h]
\centering
\caption{Comparison of model variants on LineMOD.}
\label{tab:main_results}
\begin{tabular}{lcc}
\toprule
Model & ADD-2cm (\%) & Mean ADD (mm) \\
\midrule
RGB (baseline) & 25.2 & 68.4 \\
RGB-Geometric & \textbf{33.1} & \textbf{52.3} \\
\midrule
RGBD-Geometric (bonus) & 38.5 & 45.1 \\
\bottomrule
\end{tabular}
\end{table}

The RGB-Geometric model significantly outperforms the baseline RGB model, demonstrating the benefits of geometric constraints. The RGBD-Geometric model further improves performance by leveraging depth information.

\subsection{Per-Object Analysis}

\cref{fig:per_object} shows per-object ADD-2cm accuracy. Performance varies significantly across objects, with some achieving over 50\% accuracy while others remain challenging.

% Placeholder for figure
\begin{figure}[h]
\centering
\fbox{\parbox{0.9\linewidth}{\centering [Per-Object Results Figure]\\ADD-2cm accuracy for each of 13 objects}}
\caption{Per-object ADD-2cm accuracy on LineMOD. Objects with distinctive textures and shapes achieve higher accuracy.}
\label{fig:per_object}
\end{figure}

\subsection{Training Dynamics}

\cref{fig:training_curves} shows the training curves over 150 epochs. The model converges smoothly with the cosine annealing schedule.

\begin{figure}[h]
\centering
\fbox{\parbox{0.9\linewidth}{\centering [Training Curves Figure]\\Train/Val Loss, ADD Error, Accuracy}}
\caption{Training curves showing loss, ADD error, and ADD-2cm accuracy over 150 epochs.}
\label{fig:training_curves}
\end{figure}

\subsection{Augmentation Experiments}

\cref{tab:augmentation} shows the effect of data augmentation.

\begin{table}[h]
\centering
\caption{Effect of data augmentation on RGB-Geometric model.}
\label{tab:augmentation}
\begin{tabular}{lc}
\toprule
Configuration & ADD-2cm (\%) \\
\midrule
No augmentation & \textbf{33.1} \\
+ Color Jitter & 27.0 \\
+ BBox Jitter & 32.8 \\
\bottomrule
\end{tabular}
\end{table}

Interestingly, augmentation does not improve performance on LineMOD. Color jitter actually hurts accuracy, likely because the test images have similar lighting to training. Bounding box jitter maintains similar performance but does not provide gains.

\subsection{Loss Landscape Analysis}

Following Li \etal~\cite{visualloss}, we visualize the loss landscape around our trained model using filter-normalized random directions. \cref{fig:loss_landscape} shows the resulting landscape.

\begin{figure}[h]
\centering
\fbox{\parbox{0.9\linewidth}{\centering [Loss Landscape Figure]\\3D surface and 2D contour plots}}
\caption{Loss landscape visualization using filter-normalized random directions. The smooth, bowl-shaped landscape with a flat minimum indicates good convergence and generalization potential.}
\label{fig:loss_landscape}
\end{figure}

The landscape shows a smooth, bowl-shaped surface with a relatively flat minimum region. This indicates:
\begin{itemize}
    \item \textbf{Good convergence:} The trained model found a genuine local minimum.
    \item \textbf{Flat minimum:} Associated with better generalization.
    \item \textbf{Smooth landscape:} Easy optimization without chaotic oscillations.
\end{itemize}

\subsection{Qualitative Results}

\cref{fig:qualitative} shows qualitative results on test images.

\begin{figure}[h]
\centering
\fbox{\parbox{0.9\linewidth}{\centering [Qualitative Results Figure]\\Predicted (cyan) vs GT (green) 3D boxes}}
\caption{Qualitative results showing predicted pose (cyan) versus ground truth (green). Our model accurately predicts both rotation and translation.}
\label{fig:qualitative}
\end{figure}

\subsection{Error Analysis}

Analyzing the translation error breakdown reveals that Z (depth) error dominates:

\begin{table}[h]
\centering
\caption{Translation error breakdown (mm).}
\label{tab:error_breakdown}
\begin{tabular}{lccc}
\toprule
 & X Error & Y Error & Z Error \\
\midrule
Mean & 5.2 & 4.8 & 25.3 \\
\bottomrule
\end{tabular}
\end{table}

The X and Y errors are small because they are geometrically derived from accurate bounding box detections. The Z error dominates because depth estimation from monocular images remains inherently ambiguous.
